
\documentclass{article}
\usepackage{colortbl}
\usepackage{makecell}
\usepackage{multirow}
\usepackage{supertabular}

\begin{document}

\newcounter{utterance}

\twocolumn

{ \footnotesize  \setcounter{utterance}{1}
\setlength{\tabcolsep}{0pt}
\begin{supertabular}{c@{$\;$}|p{.15\linewidth}@{}p{.15\linewidth}p{.15\linewidth}p{.15\linewidth}p{.15\linewidth}p{.15\linewidth}}

    \# & $\;$A & \multicolumn{4}{c}{Game Master} & $\;\:$B\\
    \hline 

    \theutterance \stepcounter{utterance}  

    & & \multicolumn{4}{p{0.6\linewidth}}{\cellcolor[rgb]{0.9,0.9,0.9}{%
	\makecell[{{p{\linewidth}}}]{% 
	  \tt {\tiny [A$\langle$GM]}  
	 State giocando a un gioco di negoziazione in cui dovete accordarvi su come dividere una serie di oggetti tra voi e un altro giocatore.\\ \tt \\ \tt Regole:\\ \tt (a) A Lei e all'altro giocatore viene dato un insieme di oggetti. Ognuno di voi riceve anche una funzione di valore segreta, che rappresenta il valore di ciascun tipo di oggetto.\\ \tt (b) Si scambiano messaggi con l'altro giocatore per concordare chi si aggiudica gli oggetti. Potete inviare un massimo di 5 messaggi ciascuno, oppure terminare in anticipo facendo una proposta segreta in qualsiasi momento.\\ \tt (c) A ciascuno di voi viene chiesto di inviare una proposta segreta indicando gli oggetti che desiderate, formattata tra parentesi quadre come segue: "[Proposta: <numero> <nome oggetto>, <numero> <nome oggetto>, <...>]".\\ \tt (d) Se le vostre proposte sono complementari, cioè ci sono abbastanza oggetti per soddisfare entrambe le proposte, a ciascun giocatore viene assegnato un punteggio basato sulla somma dei valori degli oggetti ricevuti. In caso contrario, entrambi ricevono zero punti.\\ \tt (e) L'obiettivo è massimizzare il proprio punteggio.\\ \tt \\ \tt Cominciamo.\\ \tt \\ \tt L'insieme degli oggetti disponibili è:\\ \tt \\ \tt 2 dadi, 1 palla, 1 colla, 1 batteria\\ \tt \\ \tt La funzione valore segreta è:\\ \tt \\ \tt dado: 1, palla: 7, colla: 0, batteria: 1\\ \tt \\ \tt IMPORTANTE: i vostri messaggi, a meno che non si tratti di una proposta segreta, vengono trasmessi direttamente all'altro giocatore, quindi non includete alcuna risposta alle regole o testo di annuncio del vostro messaggio. Per fare una proposta segreta, utilizzate il formato indicato. Non utilizzare delle parentesi quadre quando si comunica all'altro giocatore, altrimenti verrà interpretata come una proposta segreta.\\ \tt  
	  } 
	   } 
	   } 
	 & \\ 
 

    \theutterance \stepcounter{utterance}  

    & \multicolumn{4}{p{0.6\linewidth}}{\cellcolor[rgb]{0.8,1,0.9}{%
	\makecell[{{p{\linewidth}}}]{% 
	  \tt {\tiny [A$\rangle$GM]}  
	 Sembra che ci siano oggetti di valore diverso, alcuni più preziosi di altri. Mi chiedo quale sarebbe una buona strategia per iniziare. Forse potremmo discutere quali oggetti sono più importanti per noi e cercare di trovare un equilibrio? Qual è il tuo approccio iniziale? 
	  } 
	   } 
	   } 
	 & & \\ 
 

    \theutterance \stepcounter{utterance}  

    & & & \multicolumn{2}{p{0.3\linewidth}}{\cellcolor[rgb]{0.95,0.95,0.95}{%
	\makecell[{{p{\linewidth}}}]{% 
	  \tt {\tiny [GM$|$GM]}  
	 continue 
	  } 
	   } 
	   } 
	 & & \\ 
 

    \theutterance \stepcounter{utterance}  

    & & & \multicolumn{2}{p{0.3\linewidth}}{\cellcolor[rgb]{0.95,0.95,0.95}{%
	\makecell[{{p{\linewidth}}}]{% 
	  \tt {\tiny [GM$|$GM]}  
	 Sembra che ci siano oggetti di valore diverso, alcuni più preziosi di altri. Mi chiedo quale sarebbe una buona strategia per iniziare. Forse potremmo discutere quali oggetti sono più importanti per noi e cercare di trovare un equilibrio? Qual è il tuo approccio iniziale? 
	  } 
	   } 
	   } 
	 & & \\ 
 

    \theutterance \stepcounter{utterance}  

    & & \multicolumn{4}{p{0.6\linewidth}}{\cellcolor[rgb]{0.9,0.9,0.9}{%
	\makecell[{{p{\linewidth}}}]{% 
	  \tt {\tiny [GM$\rangle$B]}  
	 State giocando a un gioco di negoziazione in cui dovete accordarvi su come dividere una serie di oggetti tra voi e un altro giocatore.\\ \tt \\ \tt Regole:\\ \tt (a) A Lei e all'altro giocatore viene dato un insieme di oggetti. Ognuno di voi riceve anche una funzione di valore segreta, che rappresenta il valore di ciascun tipo di oggetto.\\ \tt (b) Si scambiano messaggi con l'altro giocatore per concordare chi si aggiudica gli oggetti. Potete inviare un massimo di 5 messaggi ciascuno, oppure terminare in anticipo facendo una proposta segreta in qualsiasi momento.\\ \tt (c) A ciascuno di voi viene chiesto di inviare una proposta segreta indicando gli oggetti che desiderate, formattata tra parentesi quadre come segue: "[Proposta: <numero> <nome oggetto>, <numero> <nome oggetto>, <...>]".\\ \tt (d) Se le vostre proposte sono complementari, cioè ci sono abbastanza oggetti per soddisfare entrambe le proposte, a ciascun giocatore viene assegnato un punteggio basato sulla somma dei valori degli oggetti ricevuti. In caso contrario, entrambi ricevono zero punti.\\ \tt (e) L'obiettivo è massimizzare il proprio punteggio.\\ \tt \\ \tt Cominciamo.\\ \tt \\ \tt L'insieme degli oggetti disponibili è:\\ \tt \\ \tt 2 dadi, 1 palla, 1 colla, 1 batteria\\ \tt \\ \tt La funzione valore segreta è:\\ \tt \\ \tt dado: 0, palla: 2, colla: 8, batteria: 0\\ \tt \\ \tt IMPORTANTE: i vostri messaggi, a meno che non si tratti di una proposta segreta, vengono trasmessi direttamente all'altro giocatore, quindi non includete alcuna risposta alle regole o testo di annuncio del vostro messaggio. Per fare una proposta segreta, utilizzate il formato indicato. Non utilizzare delle parentesi quadre quando si comunica all'altro giocatore, altrimenti verrà interpretata come una proposta segreta.\\ \tt \\ \tt \\ \tt Sembra che ci siano oggetti di valore diverso, alcuni più preziosi di altri. Mi chiedo quale sarebbe una buona strategia per iniziare. Forse potremmo discutere quali oggetti sono più importanti per noi e cercare di trovare un equilibrio? Qual è il tuo approccio iniziale? 
	  } 
	   } 
	   } 
	 & \\ 
 

    \theutterance \stepcounter{utterance}  

    & & & \multicolumn{4}{p{0.6\linewidth}}{\cellcolor[rgb]{1,0.85,0.72}{%
	\makecell[{{p{\linewidth}}}]{% 
	  \tt {\tiny [GM$\langle$B]}  
	 Penso che iniziare con una discussione aperta sui nostri interessi sia una buona strategia. Non conoscendo i valori segreti dell'altro giocatore, potremmo iniziare chiedendoci cosa ognuno di noiConsidera più importante. Forse potremmo parlare della palla e della colla, che potrebbero essere oggetti più desiderabili? Inoltre, potremmo sondare l'interesse dell'altro giocatore per i dadi e la batteria, che potrebbero essere meno importanti per loro. Una possibile strategia potrebbe essere quella di offrire qualcosa in cambio di ciò che desideriamo di più, in modo da trovare un accordo che soddisfi entrambe le parti. Cosa ne pensi? 
	  } 
	   } 
	   } 
	  \\ 
 

    \theutterance \stepcounter{utterance}  

    & & & \multicolumn{2}{p{0.3\linewidth}}{\cellcolor[rgb]{0.95,0.95,0.95}{%
	\makecell[{{p{\linewidth}}}]{% 
	  \tt {\tiny [GM$|$GM]}  
	 continue 
	  } 
	   } 
	   } 
	 & & \\ 
 

    \theutterance \stepcounter{utterance}  

    & & & \multicolumn{2}{p{0.3\linewidth}}{\cellcolor[rgb]{0.95,0.95,0.95}{%
	\makecell[{{p{\linewidth}}}]{% 
	  \tt {\tiny [GM$|$GM]}  
	 Penso che iniziare con una discussione aperta sui nostri interessi sia una buona strategia. Non conoscendo i valori segreti dell'altro giocatore, potremmo iniziare chiedendoci cosa ognuno di noiConsidera più importante. Forse potremmo parlare della palla e della colla, che potrebbero essere oggetti più desiderabili? Inoltre, potremmo sondare l'interesse dell'altro giocatore per i dadi e la batteria, che potrebbero essere meno importanti per loro. Una possibile strategia potrebbe essere quella di offrire qualcosa in cambio di ciò che desideriamo di più, in modo da trovare un accordo che soddisfi entrambe le parti. Cosa ne pensi? 
	  } 
	   } 
	   } 
	 & & \\ 
 

    \theutterance \stepcounter{utterance}  

    & & \multicolumn{4}{p{0.6\linewidth}}{\cellcolor[rgb]{0.9,0.9,0.9}{%
	\makecell[{{p{\linewidth}}}]{% 
	  \tt {\tiny [A$\langle$GM]}  
	 Penso che iniziare con una discussione aperta sui nostri interessi sia una buona strategia. Non conoscendo i valori segreti dell'altro giocatore, potremmo iniziare chiedendoci cosa ognuno di noiConsidera più importante. Forse potremmo parlare della palla e della colla, che potrebbero essere oggetti più desiderabili? Inoltre, potremmo sondare l'interesse dell'altro giocatore per i dadi e la batteria, che potrebbero essere meno importanti per loro. Una possibile strategia potrebbe essere quella di offrire qualcosa in cambio di ciò che desideriamo di più, in modo da trovare un accordo che soddisfi entrambe le parti. Cosa ne pensi? 
	  } 
	   } 
	   } 
	 & \\ 
 

    \theutterance \stepcounter{utterance}  

    & \multicolumn{4}{p{0.6\linewidth}}{\cellcolor[rgb]{0.8,1,0.9}{%
	\makecell[{{p{\linewidth}}}]{% 
	  \tt {\tiny [A$\rangle$GM]}  
	 Mi sembra un approccio ragionevole. La palla sembra un oggetto molto interessante, considerando che è unico nel suo genere, mentre la colla potrebbe essere meno desiderabile. I dadi e la batteria, invece, sembrano più comuni e potrebbero essere più facili da scambiare. Forse potremmo offrire uno o entrambi i dadi in cambio di qualcos'altro che ci interessa? O magari la batteria potrebbe essere un ottimo oggetto da scambiare per qualcosa di più prezioso? Cosa pensi della possibilità di scambiare uno degli oggetti meno importanti per noi con qualcosa che potrebbe essere più importante per l'altro giocatore? 
	  } 
	   } 
	   } 
	 & & \\ 
 

    \theutterance \stepcounter{utterance}  

    & & & \multicolumn{2}{p{0.3\linewidth}}{\cellcolor[rgb]{0.95,0.95,0.95}{%
	\makecell[{{p{\linewidth}}}]{% 
	  \tt {\tiny [GM$|$GM]}  
	 continue 
	  } 
	   } 
	   } 
	 & & \\ 
 

    \theutterance \stepcounter{utterance}  

    & & & \multicolumn{2}{p{0.3\linewidth}}{\cellcolor[rgb]{0.95,0.95,0.95}{%
	\makecell[{{p{\linewidth}}}]{% 
	  \tt {\tiny [GM$|$GM]}  
	 Mi sembra un approccio ragionevole. La palla sembra un oggetto molto interessante, considerando che è unico nel suo genere, mentre la colla potrebbe essere meno desiderabile. I dadi e la batteria, invece, sembrano più comuni e potrebbero essere più facili da scambiare. Forse potremmo offrire uno o entrambi i dadi in cambio di qualcos'altro che ci interessa? O magari la batteria potrebbe essere un ottimo oggetto da scambiare per qualcosa di più prezioso? Cosa pensi della possibilità di scambiare uno degli oggetti meno importanti per noi con qualcosa che potrebbe essere più importante per l'altro giocatore? 
	  } 
	   } 
	   } 
	 & & \\ 
 

    \theutterance \stepcounter{utterance}  

    & & \multicolumn{4}{p{0.6\linewidth}}{\cellcolor[rgb]{0.9,0.9,0.9}{%
	\makecell[{{p{\linewidth}}}]{% 
	  \tt {\tiny [GM$\rangle$B]}  
	 Mi sembra un approccio ragionevole. La palla sembra un oggetto molto interessante, considerando che è unico nel suo genere, mentre la colla potrebbe essere meno desiderabile. I dadi e la batteria, invece, sembrano più comuni e potrebbero essere più facili da scambiare. Forse potremmo offrire uno o entrambi i dadi in cambio di qualcos'altro che ci interessa? O magari la batteria potrebbe essere un ottimo oggetto da scambiare per qualcosa di più prezioso? Cosa pensi della possibilità di scambiare uno degli oggetti meno importanti per noi con qualcosa che potrebbe essere più importante per l'altro giocatore? 
	  } 
	   } 
	   } 
	 & \\ 
 

    \theutterance \stepcounter{utterance}  

    & & & \multicolumn{4}{p{0.6\linewidth}}{\cellcolor[rgb]{1,0.85,0.72}{%
	\makecell[{{p{\linewidth}}}]{% 
	  \tt {\tiny [GM$\langle$B]}  
	 Sono d'accordo che la palla sia un oggetto molto interessante e la colla potrebbe essere meno desiderabile per l'altro giocatore. Offrire uno o entrambi i dadi in cambio di qualcos'altro che ci interessa potrebbe essere una buona strategia, poiché potrebbero non essere molto importanti per l'altro giocatore. La batteria potrebbe essere un buon oggetto da scambiare, poiché potrebbe essere considerata di poco valore dall'altro giocatore, ma potremmo usarla come merce di scambio per ottenere qualcosa di più prezioso. Potremmo anche provare a chiedere all'altro giocatore cosa pensa della palla e della colla, per capire se sono disposti a scambiare uno di questi oggetti con qualcos'altro. Magari potremmo trovare un accordo che ci permetta di ottenere la palla, che sembra essere l'oggetto più prezioso, in cambio di qualcosa di meno importante per noi. Cosa ne pensi di fare una proposta di scambio, tipo offrire un dado e la batteria in cambio della palla? 
	  } 
	   } 
	   } 
	  \\ 
 

    \theutterance \stepcounter{utterance}  

    & & & \multicolumn{2}{p{0.3\linewidth}}{\cellcolor[rgb]{0.95,0.95,0.95}{%
	\makecell[{{p{\linewidth}}}]{% 
	  \tt {\tiny [GM$|$GM]}  
	 continue 
	  } 
	   } 
	   } 
	 & & \\ 
 

    \theutterance \stepcounter{utterance}  

    & & & \multicolumn{2}{p{0.3\linewidth}}{\cellcolor[rgb]{0.95,0.95,0.95}{%
	\makecell[{{p{\linewidth}}}]{% 
	  \tt {\tiny [GM$|$GM]}  
	 Sono d'accordo che la palla sia un oggetto molto interessante e la colla potrebbe essere meno desiderabile per l'altro giocatore. Offrire uno o entrambi i dadi in cambio di qualcos'altro che ci interessa potrebbe essere una buona strategia, poiché potrebbero non essere molto importanti per l'altro giocatore. La batteria potrebbe essere un buon oggetto da scambiare, poiché potrebbe essere considerata di poco valore dall'altro giocatore, ma potremmo usarla come merce di scambio per ottenere qualcosa di più prezioso. Potremmo anche provare a chiedere all'altro giocatore cosa pensa della palla e della colla, per capire se sono disposti a scambiare uno di questi oggetti con qualcos'altro. Magari potremmo trovare un accordo che ci permetta di ottenere la palla, che sembra essere l'oggetto più prezioso, in cambio di qualcosa di meno importante per noi. Cosa ne pensi di fare una proposta di scambio, tipo offrire un dado e la batteria in cambio della palla? 
	  } 
	   } 
	   } 
	 & & \\ 
 

    \theutterance \stepcounter{utterance}  

    & & \multicolumn{4}{p{0.6\linewidth}}{\cellcolor[rgb]{0.9,0.9,0.9}{%
	\makecell[{{p{\linewidth}}}]{% 
	  \tt {\tiny [A$\langle$GM]}  
	 Sono d'accordo che la palla sia un oggetto molto interessante e la colla potrebbe essere meno desiderabile per l'altro giocatore. Offrire uno o entrambi i dadi in cambio di qualcos'altro che ci interessa potrebbe essere una buona strategia, poiché potrebbero non essere molto importanti per l'altro giocatore. La batteria potrebbe essere un buon oggetto da scambiare, poiché potrebbe essere considerata di poco valore dall'altro giocatore, ma potremmo usarla come merce di scambio per ottenere qualcosa di più prezioso. Potremmo anche provare a chiedere all'altro giocatore cosa pensa della palla e della colla, per capire se sono disposti a scambiare uno di questi oggetti con qualcos'altro. Magari potremmo trovare un accordo che ci permetta di ottenere la palla, che sembra essere l'oggetto più prezioso, in cambio di qualcosa di meno importante per noi. Cosa ne pensi di fare una proposta di scambio, tipo offrire un dado e la batteria in cambio della palla? 
	  } 
	   } 
	   } 
	 & \\ 
 

    \theutterance \stepcounter{utterance}  

    & \multicolumn{4}{p{0.6\linewidth}}{\cellcolor[rgb]{0.8,1,0.9}{%
	\makecell[{{p{\linewidth}}}]{% 
	  \tt {\tiny [A$\rangle$GM]}  
	 Una proposta di scambio del genere potrebbe essere un buon punto di partenza. Offrire un dado e la batteria in cambio della palla sembra una proposta ragionevole, considerando che la palla sembra essere l'oggetto più prezioso. Tuttavia, potremmo anche considerare di tenere un po' di flessibilità nella nostra proposta, nel caso in cui l'altro giocatore non sia disposto a scambiare la palla con quegli oggetti. Forse potremmo aggiungere una opzione di scambio alternativa, come offrire entrambi i dadi e la batteria in cambio della palla e della colla. In questo modo, mostriamo che siamo disposti a negoziare e a trovare un accordo che soddisfi entrambe le parti. Cosa ne pensi di provare a fare una proposta di scambio e vedere come reagisce l'altro giocatore? [Proposta: 1 palla] 
	  } 
	   } 
	   } 
	 & & \\ 
 

    \theutterance \stepcounter{utterance}  

    & & & \multicolumn{2}{p{0.3\linewidth}}{\cellcolor[rgb]{0.95,0.95,0.95}{%
	\makecell[{{p{\linewidth}}}]{% 
	  \tt {\tiny [GM$|$GM]}  
	 proposal 
	  } 
	   } 
	   } 
	 & & \\ 
 

    \theutterance \stepcounter{utterance}  

    & & & \multicolumn{2}{p{0.3\linewidth}}{\cellcolor[rgb]{0.95,0.95,0.95}{%
	\makecell[{{p{\linewidth}}}]{% 
	  \tt {\tiny [GM$|$GM]}  
	 [0, 1, 0, 0] 
	  } 
	   } 
	   } 
	 & & \\ 
 

    \theutterance \stepcounter{utterance}  

    & & \multicolumn{4}{p{0.6\linewidth}}{\cellcolor[rgb]{0.9,0.9,0.9}{%
	\makecell[{{p{\linewidth}}}]{% 
	  \tt {\tiny [GM$\rangle$B]}  
	 L'altro giocatore ha presentato una proposta segreta. Ora tocca a Lei presentare una sua proposta.\\ \tt  
	  } 
	   } 
	   } 
	 & \\ 
 

    \theutterance \stepcounter{utterance}  

    & & & \multicolumn{4}{p{0.6\linewidth}}{\cellcolor[rgb]{1,0.85,0.72}{%
	\makecell[{{p{\linewidth}}}]{% 
	  \tt {\tiny [GM$\langle$B]}  
	 [Proposta: 1 palla, 1 colla] 
	  } 
	   } 
	   } 
	  \\ 
 

    \theutterance \stepcounter{utterance}  

    & & & \multicolumn{2}{p{0.3\linewidth}}{\cellcolor[rgb]{0.95,0.95,0.95}{%
	\makecell[{{p{\linewidth}}}]{% 
	  \tt {\tiny [GM$|$GM]}  
	 proposal 
	  } 
	   } 
	   } 
	 & & \\ 
 

    \theutterance \stepcounter{utterance}  

    & & & \multicolumn{2}{p{0.3\linewidth}}{\cellcolor[rgb]{0.95,0.95,0.95}{%
	\makecell[{{p{\linewidth}}}]{% 
	  \tt {\tiny [GM$|$GM]}  
	 [0, 1, 1, 0] 
	  } 
	   } 
	   } 
	 & & \\ 
 

    \theutterance \stepcounter{utterance}  

    & & & \multicolumn{2}{p{0.3\linewidth}}{\cellcolor[rgb]{0.95,0.95,0.95}{%
	\makecell[{{p{\linewidth}}}]{% 
	  \tt {\tiny [GM$|$GM]}  
	 [[0, 1, 0, 0], [0, 1, 1, 0]] 
	  } 
	   } 
	   } 
	 & & \\ 
 

\end{supertabular}
}

\end{document}
